% Generated draft: requires mathematical and source review.
\documentclass[11pt]{article}
\usepackage[margin=25mm]{geometry}
\usepackage{fontspec}
\setmainfont{Latin Modern Roman}
\usepackage{amsmath,amssymb,mathtools,unicode-math}
\setmathfont{Latin Modern Math}
\usepackage{enumitem,needspace,xurl,hyperref}
\setlength{\parindent}{0pt}
\setlength{\parskip}{5pt}
\setlength{\emergencystretch}{3em}
\newcommand{\N}{\mathbb{N}}
\newcommand{\Z}{\mathbb{Z}}
\newcommand{\Q}{\mathbb{Q}}
\newcommand{\R}{\mathbb{R}}
\newcommand{\C}{\mathbb{C}}
\newcommand{\F}{\mathbb{F}}
\newcommand{\A}{\mathbb{A}}
\newcommand{\PP}{\mathbb{P}}
\newcommand{\E}{\mathbb{E}}
\newcommand{\eps}{\varepsilon}
\newcommand{\dd}{\,\mathrm{d}}
\begin{document}
\section*{UnsolvedMath ID: 701206}
\textbf{Source problem code:} catalogue entry
\subsection{Definitions and mathematical statement}
Let \(G\) be a finite \(2\)-group, meaning that \(|G|=2^n\) for some integer \(n\geq 0\). Define
\[
Z(G)=\{z\in G:zg=gz\text{ for every }g\in G\},\qquad
\Omega_1(G)=\langle g\in G:g^2=1\rangle,\qquad
G^2=\langle g^2:g\in G\rangle.
\]
Here \(G^2\) denotes the subgroup generated by squares. For a finite abelian \(2\)-group \(A\), use the convention
\[
\operatorname{rk}(A)=\dim_{\F_2}(A/A^2),
\]
which is its minimum number of generators. The supplied source does not explicitly define its rank convention; this draft uses the preceding standard convention for the two abelian groups in the question. The quotient \(G/G^2\) is elementary abelian.

The question attributed to Ya. G. Berkovich is whether every finite \(2\)-group \(G\) satisfying \(\Omega_1(G)\leq Z(G)\) satisfies
\[
\operatorname{rk}(G/G^2)\leq 2\operatorname{rk}(Z(G)),
\]
or, equivalently,
\[
\dim_{\F_2}(G/G^2)\leq 2\dim_{\F_2}\bigl(Z(G)/Z(G)^2\bigr).
\]
The supplied registry labels this problem open; that label is not an independently verified assessment of its current status.
\subsection{Short English statement}
Must the minimum number of generators of a finite \(2\)-group whose involutions are all central be at most twice the minimum number of generators of its center? This uses the rank convention specified above.
\subsection{Sources}
\begin{itemize}
\item Supplied local UnsolvedMath registry record: permanent ID \(701206\), exact source problem code \texttt{catalogue entry}, title \emph{Kourovka Notebook Problem 16.14}; upstream revision \texttt{d371ee917ac9647fbb06cd6fc529deec6f3dda82}. The record was created and updated on \(2026\text{-}10\text{-}16\) and labels the problem open.
\item \emph{The Kourovka Notebook (Unsolved Problems in Group Theory)}, Problem 16.14, proposed by Ya. G. Berkovich, 16th issue (2006). The supplied background also identifies the 21st issue edition (2026), arXiv:1401.0300v48. These bibliographic details are supplied attribution and have not been independently verified.
\item The supplied collection description attributes editing of Issues 1--21 (1965--2026) to E. I. Khukhro and V. D. Mazurov. No source excerpt beyond the supplied record, explicit reference URL, or dated literature-verification assessment is provided.
\end{itemize}
\end{document}
